\documentclass[12pt,a4paper]{report}

\usepackage{iftex}
\ifPDFTeX
  \usepackage[utf8]{inputenc}
  \usepackage[T1]{fontenc}
  \usepackage{newtxtext}
  \usepackage{newtxmath}
\else
  \usepackage{fontspec}
  \setmainfont{Times New Roman}
  \setmonofont{Courier New}
\fi
\usepackage{geometry}
\usepackage{xcolor}
\usepackage{textcomp}
\usepackage{listings}
\usepackage{caption}
\usepackage{hyperref}
\usepackage{array}
\usepackage{booktabs}
\usepackage{setspace}

\geometry{left=4cm,right=3cm,top=3cm,bottom=3cm}
\onehalfspacing
\emergencystretch=3em

\makeatletter
\renewcommand{\@makechapterhead}[1]{%
  \vspace*{10pt}%
  {\parindent \z@ \centering \normalfont
    \bfseries\fontsize{14}{16}\selectfont
    \chaptername\ \thechapter\par
    \vspace{6pt}%
    #1\par\nobreak
    \vspace{18pt}%
  }}
\renewcommand{\section}{\@startsection{section}{1}{\z@}%
  {-2.5ex \@plus -1ex \@minus -.2ex}%
  {1.3ex \@plus .2ex}%
  {\normalfont\bfseries\fontsize{12}{14}\selectfont}}
\renewcommand{\subsection}{\@startsection{subsection}{2}{\z@}%
  {-2.25ex\@plus -1ex \@minus -.2ex}%
  {1ex \@plus .2ex}%
  {\normalfont\bfseries\fontsize{12}{14}\selectfont}}
\makeatother

\definecolor{codebg}{RGB}{248,248,248}
\definecolor{codeborder}{RGB}{40,40,40}
\definecolor{commentgreen}{RGB}{0,110,70}
\definecolor{keywordblue}{RGB}{0,70,150}
\definecolor{stringred}{RGB}{150,40,40}

\lstdefinelanguage{Dart}{
  morekeywords={
    abstract,as,assert,async,await,bool,break,case,catch,class,const,continue,
    default,deferred,do,double,dynamic,else,enum,export,extends,extension,
    external,factory,false,final,finally,for,Future,get,hide,if,implements,
    import,in,int,interface,is,late,library,List,Map,mixin,new,null,on,operator,
    part,required,return,set,show,static,String,super,switch,sync,this,throw,
    true,try,typedef,var,void,while,with,yield
  },
  sensitive=true,
  morecomment=[l]{//},
  morecomment=[s]{/*}{*/},
  morestring=[b]',
  morestring=[b]"
}

\lstset{
  basicstyle=\ttfamily\footnotesize,
  backgroundcolor=\color{codebg},
  frame=single,
  rulecolor=\color{codeborder},
  breaklines=true,
  breakatwhitespace=false,
  columns=fullflexible,
  keepspaces=true,
  showstringspaces=false,
  tabsize=2,
  numbers=none,
  keywordstyle=\color{keywordblue}\bfseries,
  commentstyle=\color{commentgreen},
  stringstyle=\color{stringred},
  captionpos=b
}

\hypersetup{
  colorlinks=true,
  linkcolor=black,
  urlcolor=blue
}

\begin{document}

\renewcommand{\chaptername}{BAB}
\setcounter{chapter}{3}
\chapter{Implementasi dan Pengujian Sistem}

\section{Gambaran Umum Implementasi}

Bab ini menjelaskan implementasi sistem pemantauan kesehatan berbasis Internet of Things
(IoT) yang dikembangkan pada proyek ini. Sistem terdiri dari perangkat keras berbasis
ESP32, sensor Galvanic Skin Response (GSR), sensor Electromyography (EMG), layar TFT
ST7735, Firebase Realtime Database, serta aplikasi mobile berbasis Flutter. Implementasi
dibuat mengikuti rancangan pada Bab II dan Bab III, yaitu pembacaan sinyal GSR dan EMG,
kompresi menggunakan metode \textit{Compressive Sensing} (CS), pengiriman data ke
Firebase, rekonstruksi sinyal pada aplikasi mobile, serta visualisasi data secara
\textit{real-time}.

Alur kerja sistem yang diimplementasikan adalah sebagai berikut.

\begin{enumerate}
  \item ESP32 melakukan inisialisasi pin sensor, layar TFT, Wi-Fi, Firebase, dan parameter CS.
  \item Sensor GSR dan EMG dibaca melalui ADC 12-bit pada ESP32.
  \item Data sensor difilter dan divalidasi agar pin kosong atau sinyal tidak stabil tidak dikirim sebagai data valid.
  \item Window data berukuran 64 sampel diubah ke domain wavelet menggunakan DWT Haar.
  \item Koefisien DWT dikompresi menggunakan matriks acak Gaussian menjadi 24 nilai terkompresi.
  \item Paket JSON dikirim ke Firebase Realtime Database.
  \item Aplikasi Flutter membaca data Firebase, menjalankan rekonstruksi OMP, lalu menampilkan nilai GSR, EMG, grafik, klasifikasi, dan laporan PDF.
\end{enumerate}

\section{Implementasi Perangkat Keras ESP32}

\subsection{Konfigurasi Library dan Pin}

Program ESP32 dikembangkan menggunakan PlatformIO dengan framework Arduino. Firmware
menggunakan pustaka Firebase ESP Client untuk komunikasi ke Firebase, Adafruit GFX dan
Adafruit ST7735 untuk tampilan TFT, serta SPI dan Wi-Fi bawaan ESP32. Sensor GSR
dihubungkan ke GPIO34 dan sensor EMG ke GPIO35. Kedua pin tersebut merupakan pin ADC
input-only pada ESP32, sehingga pada rangkaian fisik diperlukan \textit{pull-down}
eksternal agar pin tidak mengambang saat sensor belum tersambung.

\lstinputlisting[
  language=C++,
  firstline=1,
  lastline=23,
  caption={Konfigurasi library, pin TFT, pin sensor, dan parameter validasi ADC pada ESP32.}
]{../esp32/i_care_esp32/i_care_esp32.ino}

Pada potongan program tersebut, GPIO34 digunakan sebagai input sensor GSR dan GPIO35
sebagai input sensor EMG. Variabel \texttt{ADC\_FILTER\_SAMPLES} digunakan untuk proses
median filter, sedangkan \texttt{GSR\_PRESENT\_ADC\_MIN} dan
\texttt{EMG\_PRESENT\_ADC\_MIN} digunakan sebagai ambang validasi keberadaan sinyal
sensor. Mekanisme validasi ini penting karena pin analog ESP32 dapat menghasilkan nilai
acak ketika jalur sinyal tidak diberi referensi ke GND.

\subsection{Konfigurasi Compressive Sensing}

Konfigurasi CS berada pada file \texttt{cs\_config.h}. Parameter utama yang digunakan
adalah \texttt{CS\_N = 64} sebagai jumlah sampel asli dalam satu window dan
\texttt{CS\_M = 24} sebagai jumlah sampel terkompresi. Dengan demikian, setiap sensor
mengirim 24 nilai hasil kompresi dari 64 sampel awal.

\lstinputlisting[
  language=C++,
  firstline=4,
  lastline=20,
  caption={Parameter utama Compressive Sensing pada firmware ESP32.}
]{../esp32/i_care_esp32/cs_config.h}

Nilai rasio kompresi yang digunakan dapat dihitung sebagai berikut.

\[
CR = \frac{N}{M} = \frac{64}{24} = 2{,}67
\]

Artinya, data yang dikirimkan lebih kecil dibandingkan data mentah, sehingga sesuai
dengan tujuan perancangan sistem untuk menekan beban transmisi data pada perangkat IoT.

\subsection{Transformasi DWT dan Proyeksi Gaussian}

Sesuai rancangan sistem, sinyal tidak langsung dikirim dalam bentuk mentah. Firmware
terlebih dahulu membentuk basis wavelet Haar, mengubah data sensor ke koefisien DWT,
kemudian memproyeksikan koefisien tersebut menggunakan matriks acak Gaussian. Hasil
proyeksi inilah yang dikirim sebagai paket terkompresi.

\lstinputlisting[
  language=C++,
  firstline=44,
  lastline=64,
  caption={Pembentukan matriks Gaussian, basis DWT Haar, dan fungsi kompresi sinyal.}
]{../esp32/i_care_esp32/cs_config.h}

\lstinputlisting[
  language=C++,
  firstline=86,
  lastline=104,
  caption={Transformasi DWT dan proyeksi Gaussian untuk menghasilkan data terkompresi.}
]{../esp32/i_care_esp32/cs_config.h}

Fungsi \texttt{cs\_generate\_measurement\_matrix()} menghasilkan matriks Gaussian
\(\Phi\), sedangkan \texttt{cs\_generate\_wavelet\_basis()} membentuk matriks basis
DWT Haar \(\Psi\). Pada fungsi \texttt{cs\_compress\_signal()}, data sensor \(x\)
diubah menjadi koefisien DWT \(s\), kemudian dikompresi menjadi \(y\) menggunakan
perkalian matriks Gaussian. Secara konseptual, proses ini dapat dituliskan sebagai:

\[
s = \Psi^T x
\]

\[
y = \Phi s
\]

Dengan alur tersebut, firmware ESP32 menjalankan bagian kompresi pada sisi perangkat
(\textit{edge}), sedangkan proses rekonstruksi dilakukan pada aplikasi mobile.

\subsection{Pembacaan, Filter, dan Validasi Sensor}

Pembacaan sensor dilakukan secara periodik dalam fungsi \texttt{fillCsBuffer()}. Setiap
nilai ADC dibaca menggunakan median filter untuk mengurangi lonjakan sesaat. Setelah itu,
nilai ADC dipetakan ke satuan GSR dalam \textmu S dan EMG dalam \textmu V. Sistem juga
menghitung nilai minimum, maksimum, rata-rata ADC, dan jumlah sampel yang berada pada
kondisi rail untuk memastikan sinyal sensor valid.

\lstinputlisting[
  language=C++,
  firstline=400,
  lastline=438,
  caption={Akuisisi, filter, pemetaan nilai sensor, dan validasi sinyal GSR serta EMG.}
]{../esp32/i_care_esp32/i_care_esp32.ino}

\lstinputlisting[
  language=C++,
  firstline=443,
  lastline=480,
  caption={Validasi status sensor dan pembersihan buffer ketika sinyal tidak valid.}
]{../esp32/i_care_esp32/i_care_esp32.ino}

Jika sinyal GSR atau EMG tidak valid, buffer sensor tersebut diisi dengan nilai nol dan
status validasi dikirim ke Firebase. Dengan demikian, aplikasi tidak menampilkan nilai
sensor palsu saat sensor belum tersambung atau sinyal tidak stabil.

\subsection{Pengiriman Data ke Firebase}

Setelah proses akuisisi dan kompresi selesai, ESP32 mengirim data ke Firebase dalam
format JSON. Data dikirim ke node \texttt{health\_monitoring/latest} untuk pembacaan
terbaru dan ke node history perangkat untuk penyimpanan riwayat.

\lstinputlisting[
  language=C++,
  firstline=548,
  lastline=590,
  caption={Pembuatan payload JSON dan pengiriman data terkompresi ke Firebase.}
]{../esp32/i_care_esp32/i_care_esp32.ino}

Payload yang dikirim berisi metadata CS seperti \texttt{cs\_n}, \texttt{cs\_m},
\texttt{cs\_seed}, \texttt{cs\_transform}, dan \texttt{cs\_measurement}. Selain itu,
payload juga menyertakan array \texttt{y\_gsr} dan \texttt{y\_emg} sebagai hasil
kompresi, nilai ringkasan sensor dari ESP32, status validasi sinyal, timestamp, ID
perangkat, RSSI, dan waktu aktif perangkat.

\subsection{Inisialisasi dan Loop Utama ESP32}

Fungsi \texttt{setup()} bertugas menginisialisasi TFT, ADC, Wi-Fi, Firebase, basis DWT,
dan matriks Gaussian. Fungsi \texttt{loop()} menjalankan satu siklus pemantauan yang
terdiri dari pembacaan sensor, kompresi, pengiriman data ke Firebase, pencetakan data
diagnostik pada Serial Monitor, dan pembaruan tampilan TFT.

\lstinputlisting[
  language=C++,
  firstline=658,
  lastline=704,
  caption={Inisialisasi perangkat dan siklus utama pemantauan pada ESP32.}
]{../esp32/i_care_esp32/i_care_esp32.ino}

\lstinputlisting[
  language=C++,
  firstline=707,
  lastline=731,
  caption={Siklus utama pembacaan sensor, kompresi, dan persiapan upload data.}
]{../esp32/i_care_esp32/i_care_esp32.ino}

Jika Wi-Fi atau Firebase belum siap, sistem tetap menjalankan pembacaan sensor lokal dan
menampilkan data pada TFT. Hal ini membuat perangkat tetap dapat digunakan untuk monitoring
lokal walaupun koneksi internet sedang terputus.

\section{Implementasi Firebase Realtime Database}

Firebase Realtime Database digunakan sebagai penghubung antara ESP32 dan aplikasi mobile.
Aplikasi membaca node \texttt{health\_monitoring/latest} secara real-time. Node ini
diperbarui oleh ESP32 setiap kali satu window data selesai dikompresi dan dikirim.

\lstinputlisting[
  language=Dart,
  firstline=20,
  lastline=78,
  caption={Data source Flutter untuk membaca data dari Firebase Realtime Database.}
]{../lib/data/datasources/firebase_health_remote_data_source.dart}

Pada bagian tersebut, \texttt{FirebaseHealthRemoteDataSource} membuat koneksi ke Firebase
Realtime Database dan melakukan \textit{streaming} data menggunakan \texttt{onValue}.
Jika payload memiliki \texttt{cs\_enabled = true}, aplikasi melanjutkan data tersebut ke
proses rekonstruksi CS.

\section{Implementasi Rekonstruksi Sinyal pada Aplikasi Mobile}

\subsection{Rekonstruksi OMP dan Inverse DWT}

Aplikasi mobile bertugas melakukan rekonstruksi sinyal dari data terkompresi. Data
\texttt{y\_gsr} dan \texttt{y\_emg} yang diterima dari Firebase diproses menggunakan
algoritma Orthogonal Matching Pursuit (OMP). Hasil OMP berupa koefisien sparse pada
domain wavelet, kemudian dikembalikan ke domain waktu menggunakan inverse DWT.

\lstinputlisting[
  language=Dart,
  firstline=92,
  lastline=127,
  caption={Pemanggilan proses rekonstruksi CS pada data source aplikasi Flutter.}
]{../lib/data/datasources/firebase_health_remote_data_source.dart}

\lstinputlisting[
  language=Dart,
  firstline=185,
  lastline=205,
  caption={Rekonstruksi sinyal GSR dan EMG menggunakan OMP dan inverse DWT.}
]{../lib/data/datasources/firebase_health_remote_data_source.dart}

Nilai rata-rata GSR dan EMG yang ditampilkan pada aplikasi berasal dari sinyal hasil
rekonstruksi. Apabila status validasi sensor dari ESP32 menyatakan sinyal tidak valid,
aplikasi mengatur nilai sensor menjadi nol agar dashboard tidak menampilkan hasil yang
tidak dapat dipertanggungjawabkan.

\subsection{Model Data Sensor}

Payload Firebase diubah menjadi model aplikasi melalui \texttt{HealthMetricsModel}.
Model ini membaca nilai sensor, timestamp, status validasi, parameter CS, dan array hasil
kompresi. Model ini juga menjaga kompatibilitas dengan beberapa format key seperti
\texttt{gsr}, \texttt{gsr\_value}, dan \texttt{gsrValue}.

\lstinputlisting[
  language=Dart,
  firstline=3,
  lastline=79,
  caption={Model data untuk payload sensor GSR, EMG, dan metadata CS.}
]{../lib/data/models/health_metrics_model.dart}

\section{Implementasi Dashboard Aplikasi}

Dashboard menampilkan ringkasan kondisi GSR dan EMG, nilai sensor terbaru, serta grafik
real-time. Data dashboard dikelola oleh \texttt{DashboardViewModel}. ViewModel ini
menyimpan data terbaru, riwayat GSR, riwayat EMG, status CS, rasio kompresi, dan sinyal
hasil rekonstruksi.

\lstinputlisting[
  language=Dart,
  firstline=101,
  lastline=145,
  caption={Penyimpanan data real-time dan history pada DashboardViewModel.}
]{../lib/presentation/viewmodels/dashboard_view_model.dart}

Halaman dashboard melakukan klasifikasi GSR dan EMG berdasarkan nilai terbaru yang sudah
diterima dari ViewModel. Nilai tersebut kemudian digunakan untuk menampilkan kartu status,
nilai sensor, dan grafik.

\lstinputlisting[
  language=Dart,
  firstline=47,
  lastline=84,
  caption={Penggunaan nilai GSR dan EMG pada halaman dashboard.}
]{../lib/presentation/pages/dashboard_page.dart}

\section{Implementasi Klasifikasi Stres}

Klasifikasi sensor diseragamkan pada aplikasi dan firmware. Pada aplikasi, aturan
klasifikasi berada pada \texttt{StressLevelMapper}. Klasifikasi GSR menggunakan tiga
kategori, yaitu Low, Moderate, dan High. Klasifikasi EMG menggunakan dua kategori, yaitu
Normal dan Stress.

\begin{table}[h]
\centering
\caption{Rentang klasifikasi sensor GSR dan EMG.}
\begin{tabular}{>{\raggedright\arraybackslash}p{3cm}p{4cm}p{4cm}}
\toprule
\textbf{Sensor} & \textbf{Rentang} & \textbf{Status} \\
\midrule
GSR & 1--5 \textmu S & Low \\
GSR & $>$5--12 \textmu S & Moderate \\
GSR & $>$12--20 \textmu S & High \\
EMG & $<$200 \textmu V & Normal \\
EMG & $>$200 \textmu V & Stress \\
\bottomrule
\end{tabular}
\end{table}

\lstinputlisting[
  language=Dart,
  firstline=34,
  lastline=83,
  caption={Implementasi klasifikasi GSR, EMG, dan kombinasi status stres pada aplikasi.}
]{../lib/core/utils/stress_level.dart}

\lstinputlisting[
  language=Dart,
  firstline=86,
  lastline=141,
  caption={Klasifikasi gabungan GSR dan EMG untuk interpretasi kondisi pengguna.}
]{../lib/core/utils/stress_level.dart}

Selain klasifikasi tunggal, aplikasi juga menggabungkan hasil GSR dan EMG menjadi kategori
integratif seperti \textit{Normal - Low}, \textit{Normal - Moderate},
\textit{Normal - High}, \textit{Stres - Low}, \textit{Stres - Moderate}, dan
\textit{Stres - High}. Interpretasi kategori gabungan digunakan pada halaman laporan dan
export PDF.

\section{Implementasi Laporan PDF}

Fitur laporan berada pada halaman \texttt{ReportsPage}. Aplikasi menghitung rata-rata
riwayat GSR dan EMG yang tersimpan di ViewModel, melakukan klasifikasi ulang, membuat
dokumen PDF, menyimpan file pada direktori aplikasi, lalu langsung membuka file tersebut
menggunakan \texttt{OpenFilex}.

\lstinputlisting[
  language=Dart,
  firstline=123,
  lastline=149,
  caption={Proses export PDF dan pembukaan file laporan pada aplikasi.}
]{../lib/presentation/pages/reports_page.dart}

\lstinputlisting[
  language=Dart,
  firstline=246,
  lastline=256,
  caption={Pembukaan file PDF secara langsung setelah proses export selesai.}
]{../lib/presentation/pages/reports_page.dart}

Format laporan terdiri dari header laporan kesehatan, ringkasan biometrik, nilai rata-rata
GSR, nilai rata-rata EMG, status gabungan, serta panel interpretasi kombinasi GSR dan EMG.
Dengan demikian, laporan PDF dapat digunakan sebagai ringkasan hasil pemantauan pengguna.


\end{document}
